\section{A finite-span necessary condition for a specified advantage}
\label{sec:finitespanamplitude}
The sharp quarter bound has a more restrictive form when the energy span
and the two preparation temperatures are specified.

\begin{corollary}[Finite-span cooling bound]
\label{cor:finitespanamplitude}
Under the full four-state hypotheses of
Theorem~\ref{thm:fulltripletbound}, put
\[
\Delta=E_4-E_1>0,\qquad
 d=(\beta_c-\beta_h)\Delta>0,\qquad
 p_{4,h}=\pi_4(\beta_h),
\]
for any finite $0<\beta_h<\beta_c<\beta_b$. The asymptotic
slow-amplitude advantage satisfies
\begin{equation}\label{eq:finitespanamplitude}
\mathcal A:=1-\frac{|b_2(\beta_h)|}{|b_2(\beta_c)|}
\leq p_{4,h}(1-e^{-d})
<\frac{1-e^{-d}}4
<\frac{1-e^{-\beta_b\Delta}}4.
\end{equation}
These bounds also hold when the ordering is ordinary and
$\mathcal A\leq0$.
Consequently, achieving $\mathcal A\geq g$ for a specified
$0<g<1/4$ requires
\[
(\beta_c-\beta_h)\Delta> -\log(1-4g),
\qquad p_{4,h}>g.
\]
In particular $\beta_b\Delta> -\log(1-4g)$ is necessary if the
preparation temperatures may be chosen freely on the cooling side.
These are necessary conditions only; no sufficiency or sharpness is
asserted for these finite-span bounds.
\end{corollary}
\begin{proof}
Write $P_3(\beta)=1-\pi_4(\beta)$. The proof of
Theorem~\ref{thm:fulltripletbound} established
\[
\frac{|b_2(\beta_c)|}{|b_2(\beta_h)|}
\leq\frac{P_3(\beta_c)}{P_3(\beta_h)}.
\]
Both overlaps are nonzero by the contrast theorem, so inversion gives
$\mathcal A\leq1-P_3(\beta_h)/P_3(\beta_c)$.
Shift the energy origin to $E_1=0$, leaving $E_4=\Delta$.
Then $Z_3(\beta_c)\leq Z_3(\beta_h)$ and
\[
\frac{\pi_4(\beta_c)}{P_3(\beta_c)}
=\frac{e^{-\beta_c\Delta}}{Z_3(\beta_c)}
\geq e^{-d}\frac{\pi_4(\beta_h)}{P_3(\beta_h)}.
\]
Multiplying by $P_3(\beta_h)$ gives
\[
\frac{P_3(\beta_h)}{P_3(\beta_c)}
\geq P_3(\beta_h)+e^{-d}p_{4,h},
\]
which proves the first bound. The highest-energy state has strictly
less than one-quarter Gibbs mass at finite positive $\beta_h$, so the
second inequality is strict. The last is strict because
$\beta_c-\beta_h<\beta_b$. Solving
$g<(1-e^{-d})/4$ gives the required dimensionless separation; also
$g\leq\mathcal A\leq p_{4,h}(1-e^{-d})<p_{4,h}$.
\end{proof}

For example, amplitude advantages of at least $10\%$, $20\%$, and
$24\%$ require, respectively,
\[
d>\log(5/3)\approx0.5108,\qquad
d>\log5\approx1.6094,\qquad
d>\log25\approx3.2189.
\]
They also require $p_{4,h}>0.10$, $0.20$, and $0.24$, respectively.
These are restrictions on initial inverse-temperature separation and
highest-state occupation, not guarantees that a generator satisfying
them exhibits an effect. The quantity $\mathcal A$ concerns the
asymptotic slow amplitude, not a finite-time distance or an absolute
relaxation-time reduction.
